% Компиляция: Overleaf -> Menu -> Compiler = XeLaTeX (или LuaLaTeX). Пакет minted работает в Overleaf автоматически.
\documentclass[14pt]{extarticle}

\usepackage{fontspec}
\usepackage[english,russian]{babel}
\makeatletter
\newcommand{\pickmain}{%
  \IfFontExistsTF{Times New Roman}{\setmainfont{Times New Roman}}{%
  \IfFontExistsTF{Liberation Serif}{\setmainfont{Liberation Serif}}{%
  \IfFontExistsTF{FreeSerif}{\setmainfont{FreeSerif}}{%
  \setmainfont{CMU Serif}}}}}
\newcommand{\pickmono}{%
  \IfFontExistsTF{DejaVu Sans Mono}{\setmonofont{DejaVu Sans Mono}[Scale=MatchLowercase]}{%
  \IfFontExistsTF{FreeMono}{\setmonofont{FreeMono}[Scale=MatchLowercase]}{%
  \setmonofont{CMU Typewriter Text}[Scale=MatchLowercase]}}}
\makeatother
\pickmain
\pickmono

% Поля по СТП БГУИР. Для узких полей замените строку на:
% \usepackage[a4paper,margin=12.7mm]{geometry}
\usepackage[a4paper,left=30mm,right=10mm,top=20mm,bottom=20mm]{geometry}
\usepackage{graphicx}
\usepackage{setspace}
\usepackage{indentfirst}
\usepackage{amsmath}
\usepackage{titlesec}
\usepackage{caption}
\usepackage{float}
\usepackage{tocloft}
\usepackage{minted}

\onehalfspacing
\setlength{\parindent}{1.25cm}
\titleformat{\section}{\normalsize\bfseries\raggedright}{\thesection}{1ex}{\MakeUppercase}
\titleformat{\subsection}{\normalsize\bfseries}{\thesubsection}{1ex}{}
\captionsetup{labelsep=endash,justification=centering}
\setminted{fontsize=\footnotesize,breaklines=true,frame=lines,framesep=2mm,baselinestretch=1}

\renewcommand{\contentsname}{Оглавление}
\renewcommand{\cfttoctitlefont}{\normalsize\bfseries\MakeUppercase}
\renewcommand{\cftaftertoctitle}{}
\renewcommand{\cftsecleader}{\cftdotfill{\cftdotsep}}
\renewcommand{\cftsecfont}{\normalfont}
\renewcommand{\cftsecpagefont}{\normalfont}
\setlength{\cftbeforesecskip}{2pt}

\newcommand{\screenshot}[2]{%
  \begin{figure}[H]
    \centering
    \IfFileExists{#1}{\includegraphics[width=\textwidth,height=0.42\textheight,keepaspectratio]{#1}}{\fbox{\parbox[c][5cm][c]{0.9\textwidth}{\centering Место для скриншота: #1}}}
    \caption{#2}
  \end{figure}}

\title{ММО ЛР№6}
\author{Емельяненко Е.Д.}
\date{November 2026}

\begin{document}

%======================== 1. ТИТУЛЬНЫЙ ЛИСТ ========================
\begin{titlepage}
\centering
\singlespacing
Министерство образования Республики Беларусь\\[2mm]
Учреждение образования\\
«Белорусский государственный университет информатики и радиоэлектроники»\\[8mm]
Кафедра программного обеспечения информационных технологий\\[35mm]

\textbf{ОТЧЁТ}\\
по лабораторной работе №6\\
по дисциплине «Методы машинного обучения»\\[3mm]
\textbf{Случайный лес}\\[3mm]

\begin{flushright}
Выполнил: студент гр. 351002\\
Емельяненко Е.Д.\\[5mm]
Проверил: Лебедевич А.В.
\end{flushright}

\vfill
Минск 2026
\end{titlepage}

\setcounter{page}{2}
\tableofcontents
\newpage

%======================== 2. ЦЕЛЬ ========================
\section{Цель выполнения лабораторной работы}

Изучить регрессию ансамблем деревьев и получить практические навыки её решения на табличных данных: подготовить признаки, собрать случайный лес в конвейере со стандартизацией, обучить модель с параметрами по умолчанию, оценить качество по MSE, MAE, RMSE, $R^2$ и MAPE, подобрать число деревьев, глубину и размер листа через \texttt{GridSearchCV} и разобрать вклад признаков.

%======================== 3. ТЕОРИЯ ========================
\section{Теория}

\subsection{Постановка задачи}
В регрессии каждому объекту $x_i$ сопоставлено число $y_i$, а модель выдаёт прогноз $\hat y(x)$. Качество смотрят на отложенной тестовой выборке: она не участвует ни в обучении, ни в выборе гиперпараметров. Нижняя опора "--- константный прогноз. Если всем объектам приписать среднее цели на обучении, на обучающей части $R^2 = 0$, а RMSE равен разбросу ответов вокруг этого среднего. Модель имеет смысл, пока её ошибка заметно меньше этого уровня.

\subsection{Дерево регрессии}
Дерево рекурсивно делит выборку по одному признаку и порогу. Для регрессии кандидат в разбиение оценивают взвешенной среднеквадратичной ошибкой потомков. В узле из $n$ объектов с левой и правой частями размеров $n_L$ и $n_R$ она равна
\[
\frac{n_L}{n}\,\mathrm{MSE}_L + \frac{n_R}{n}\,\mathrm{MSE}_R,
\]
где MSE части "--- средний квадрат отклонения ответов от их же среднего. Выбирают разбиение с наименьшей такой величиной. Лист прогнозирует среднее $y$ попавших в него объектов обучения.

Дерево сравнивает признак с порогом, поэтому монотонное преобразование одного столбца (логарифм, стандартизация) не меняет набор допустимых разбиений: порог просто пересчитывается в новых единицах. Глубина и минимальный размер листа ограничивают число разбиений. Без них дерево доходит до листьев из одного объекта и на обучении повторяет ответы почти точно, а на новых данных ошибка растёт.

Важность признака "--- суммарное взвешенное падение MSE по всем узлам, где признак был выбран, усреднённое по деревьям и нормированное на единицу. Большая важность означает, что признак часто даёт чистые по ответу группы. Она не показывает знак влияния и завышает один столбец, если рядом стоит его почти точная копия: разбиения достаются то одному, то другому.

\subsection{Случайный лес}
Случайный лес усредняет $B$ деревьев регрессии:
\[
\hat y(x) = \frac{1}{B}\sum_{b=1}^{B} T_b(x).
\]
Дерево $T_b$ учится на бутстреп-выборке: $N$ объектов берутся из обучения с возвращением, так что часть строк в дерево не попадает. Усреднение снижает разброс отдельных глубоких деревьев. В классической постановке Бреймана на каждом разбиении ещё рассматривают случайное подмножество признаков. В \texttt{RandomForestRegressor} из \texttt{sklearn} параметр \texttt{max\_features} по умолчанию равен $1{,}0$: в разбиение допускаются все признаки, и случайность леса даёт бутстреп. Число деревьев по умолчанию равно 100, глубина не ограничена, лист может содержать один объект.

Рост $B$ стабилизирует среднее и почти не увеличивает смещение. Ограничение глубины и увеличение минимального листа упрощают отдельные деревья: обучение запоминается хуже, а слишком жёсткое ограничение поднимает ошибку и на тесте, если зависимость не укладывается в короткие правила.

\subsection{Метрики качества}
Пусть на выборке из $n$ объектов ответы $y_i$, прогнозы $\hat y_i$, а $\bar y$ "--- среднее ответов этой же выборки. Тогда
\begin{align*}
\mathrm{MSE} &= \frac{1}{n}\sum_{i=1}^{n}(y_i - \hat y_i)^2, &
\mathrm{RMSE} &= \sqrt{\mathrm{MSE}}, \\[1mm]
\mathrm{MAE} &= \frac{1}{n}\sum_{i=1}^{n}|y_i - \hat y_i|, &
R^2 &= 1 - \frac{\sum_{i=1}^{n}(y_i - \hat y_i)^2}{\sum_{i=1}^{n}(y_i - \bar y)^2}.
\end{align*}
MSE и RMSE штрафуют крупные промахи сильнее, чем мелкие. RMSE выражен в единицах цели, MAE "--- средняя абсолютная ошибка в тех же единицах. $R^2$ показывает, какую долю разброса цели модель сняла относительно константы $\bar y$. Значение $1$ означает точное совпадение, значение $0$ "--- уровень константного прогноза.

MAPE "--- средняя абсолютная ошибка в процентах от истинного ответа:
\[
\mathrm{MAPE} = \frac{100}{n}\sum_{i=1}^{n}\left|\frac{y_i - \hat y_i}{y_i}\right|.
\]
Метрика определена, пока $y_i \neq 0$. На шкале рейтинга, где все значения положительны, это выполнено. Одинаковый промах в пунктах даёт больший MAPE у игрока с низким рейтингом.

Метрики на обучении и на тесте читают вместе. Если обучение почти безошибочно, а тест заметно хуже, деревья запомнили выборку. Если обе ошибки близки и малы относительно разброса цели, лес обобщает.

\subsection{Подготовка данных и выбор гиперпараметров}
Пропуски заполняют правилом, которое не подглядывает в тест. Медиана и мода, посчитанные по всей таблице до разбиения, уже содержат тестовые строки. Поэтому \texttt{SimpleImputer} ставят внутрь \texttt{Pipeline}: на кросс-валидации статистики считаются по обучающим фолдам, а отложенный тест не влияет на них вообще. То же относится к \texttt{StandardScaler}. Дереву масштаб не нужен, но стандартизатор входит в условие работы и живёт в том же конвейере, что и лес. Категорию с небольшим числом уровней кодируют индикаторами. Один уровень оставляют опорным (\texttt{drop='first'}), чтобы столбцы не дублировали друг друга.

Преобразование, которому не нужна статистика выборки, можно сделать до разбиения: достать год из строки даты, пометить пропуск отдельным признаком, свернуть редкие уровни в группы. Заполнять признак медианой внутри групп целевой переменной нельзя: в заполненное значение попадает сама цель.

Гиперпараметры перебирают сеткой \texttt{GridSearchCV} только по обучающей части. Для регрессии критерий "--- RMSE (\texttt{neg\_root\_mean\_squared\_error}: у скорера знак обратный, потому что поиск максимизирует оценку). Лучшую конфигурацию заново обучают на всей обучающей выборке и один раз проверяют на тесте.

%======================== 4. ЗАДАНИЕ ========================
\section{Формулировка индивидуального задания}

Вариант "--- случайный лес, задача регрессии. Данные "--- \texttt{fifa.csv}, тот же набор FIFA, что в лабораторных по разведочному анализу и градиентному спуску: 18\,207 игроков и 18 столбцов. Целевая переменная "--- общий рейтинг \texttt{Overall}.

\begin{enumerate}
\item Подготовить и проанализировать набор данных. Обработать пропуски, стандартизировать числовые признаки и закодировать категориальные. Масштабирование выполнить, если оно нужно постановке.
\item Обучить случайный лес с исходными параметрами на тренировочной части и получить прогнозы на тестовой.
\item Оценить качество. Для регрессии в задании указаны MSE, RMSE и $R^2$; дополнительно считаются MAE и MAPE.
\item Визуализировать результаты. Для регрессии вместо матрицы ошибок строятся диаграмма <<факт "--- прогноз>> и гистограмма остатков.
\item Подобрать гиперпараметры: число деревьев, глубину и размер листа.
\item Перебор выполнить через \texttt{GridSearchCV}.
\item Повторить обучение и оценку с найденными параметрами.
\item Разобрать важность признаков и построить её график. Определить, какие признаки дают основной вклад.
\end{enumerate}

Пропуски есть у \texttt{Club} (241), \texttt{Value} (252), \texttt{Contract Valid Until} (289), \texttt{International Reputation} и \texttt{Skill Moves} (по 48). Их закрывает импутатор внутри конвейера. \texttt{ID} и \texttt{Name} "--- идентификаторы. \texttt{Nationality} (164 уровня) и \texttt{Club} (651 уровень) в модель не входят: индикаторы раздули бы матрицу. Двадцать семь позиций сворачиваются в группы GK, DEF, MID и FWD. Пропуск клуба сохраняется отдельным признаком свободного агента. \texttt{Value} не заполняется медианой по группам \texttt{Overall}: рейтинг является целью. Денежные столбцы перед стандартизацией проходят \texttt{log1p}.

%======================== 5. КОД ========================
\section{Код решения индивидуального задания}

\subsection{Загрузка и обзор (пункт 1)}
\begin{minted}{python}
import numpy as np
import pandas as pd
import matplotlib.pyplot as plt

from sklearn.compose import ColumnTransformer
from sklearn.ensemble import RandomForestRegressor
from sklearn.impute import SimpleImputer
from sklearn.metrics import (
    mean_absolute_error,
    mean_absolute_percentage_error,
    mean_squared_error,
    r2_score,
)
from sklearn.model_selection import GridSearchCV, train_test_split
from sklearn.pipeline import Pipeline
from sklearn.preprocessing import FunctionTransformer, OneHotEncoder, StandardScaler

raw = pd.read_csv('fifa.csv')
print(raw.shape)
raw.info()

na = raw.isna().sum()
na = na[na > 0].to_frame('Пропуски')
na['Доля, %'] = (na['Пропуски'] / len(raw) * 100).round(2)
na

num_preview = [
    'Age', 'Potential', 'Value', 'Wage', 'International Reputation',
    'Skill Moves', 'Joined', 'Height', 'Weight', 'Release Clause',
]
corr = raw[num_preview + ['Overall']].corr(numeric_only=True)['Overall']
corr.sort_values(ascending=False).round(3)
\end{minted}

\subsection{Очистка без статистики выборки (пункт 1)}
\begin{minted}{python}
df = raw.copy()

df['Contract Year'] = pd.to_numeric(
    df['Contract Valid Until'].astype(str).str.extract(r'(\d{4})')[0],
    errors='coerce',
)
df['Is_Free_Agent'] = df['Club'].isna().astype(int)

pos_map = {
    'GK': 'GK',
    'CB': 'DEF', 'LCB': 'DEF', 'RCB': 'DEF', 'LB': 'DEF',
    'RB': 'DEF', 'LWB': 'DEF', 'RWB': 'DEF',
    'CM': 'MID', 'LCM': 'MID', 'RCM': 'MID', 'CDM': 'MID',
    'LDM': 'MID', 'RDM': 'MID', 'CAM': 'MID', 'LAM': 'MID',
    'RAM': 'MID', 'LM': 'MID', 'RM': 'MID',
    'ST': 'FWD', 'LS': 'FWD', 'RS': 'FWD', 'CF': 'FWD',
    'LF': 'FWD', 'RF': 'FWD', 'LW': 'FWD', 'RW': 'FWD',
}
df['Position_Group'] = df['Position'].map(pos_map).fillna('Other')

df = df.drop(columns=[
    'ID', 'Name', 'Nationality', 'Club', 'Position', 'Contract Valid Until',
])

print('Пропусков после подготовки:', int(df.isna().sum().sum()))
df.isna().sum()[lambda s: s > 0]
\end{minted}

\subsection{Разбиение и конвейер (пункты 1 и 2)}
\begin{minted}{python}
X = df.drop(columns=['Overall'])
y = df['Overall']

X_train, X_test, y_train, y_test = train_test_split(
    X, y, test_size=0.2, random_state=42
)

money_cols = ['Value', 'Wage', 'Release Clause']
num_cols = [
    'Age', 'Potential', 'International Reputation', 'Skill Moves',
    'Joined', 'Contract Year', 'Height', 'Weight', 'Is_Free_Agent',
]
cat_cols = ['Preferred Foot', 'Position_Group']

def make_pipeline():
    money_pipe = Pipeline([
        ('imputer', SimpleImputer(strategy='median')),
        ('log', FunctionTransformer(np.log1p, feature_names_out='one-to-one')),
        ('scaler', StandardScaler()),
    ])
    num_pipe = Pipeline([
        ('imputer', SimpleImputer(strategy='median')),
        ('scaler', StandardScaler()),
    ])
    cat_pipe = Pipeline([
        ('imputer', SimpleImputer(strategy='most_frequent')),
        ('onehot', OneHotEncoder(
            drop='first', handle_unknown='ignore', sparse_output=False,
        )),
    ])
    preprocess = ColumnTransformer([
        ('money', money_pipe, money_cols),
        ('num', num_pipe, num_cols),
        ('cat', cat_pipe, cat_cols),
    ])
    return Pipeline([
        ('prep', preprocess),
        ('rf', RandomForestRegressor(random_state=42, n_jobs=-1)),
    ])

baseline = make_pipeline()
baseline.fit(X_train, y_train)
pred_train = baseline.predict(X_train)
pred_test = baseline.predict(X_test)
\end{minted}

Медиана, \texttt{StandardScaler} и мода обучаются только на тех строках, на которых в данный момент учится лес. Для денег порядок шагов такой: медиана, \texttt{log1p}, стандартизация. Остальные числа идут без логарифма. \texttt{Preferred Foot} и группа позиции кодируются индикаторами с опорным уровнем. Перебор дальше использует тот же \texttt{make\_pipeline}, поэтому базовая модель и сетка готовят признаки одинаково.

\subsection{Метрики и графики базовой модели (пункты 3 и 4)}
\begin{minted}{python}
def regression_metrics(y_true, y_pred):
    y_true = np.asarray(y_true, dtype=float)
    y_pred = np.asarray(y_pred, dtype=float)
    mse = mean_squared_error(y_true, y_pred)
    return pd.Series({
        'MSE': mse,
        'MAE': mean_absolute_error(y_true, y_pred),
        'RMSE': float(np.sqrt(mse)),
        'R2': r2_score(y_true, y_pred),
        'MAPE, %': mean_absolute_percentage_error(y_true, y_pred) * 100,
    })

baseline_train = regression_metrics(y_train, pred_train)
baseline_test = regression_metrics(y_test, pred_test)
pd.DataFrame({'train': baseline_train, 'test': baseline_test}).round(4)

residuals = y_test.to_numpy() - pred_test
fig, axes = plt.subplots(1, 2, figsize=(12, 4))
axes[0].scatter(y_test, pred_test, alpha=0.25, s=12)
lims = [float(y_test.min()), float(y_test.max())]
axes[0].plot(lims, lims, color='black', linewidth=1)
axes[0].set_xlabel('Факт, Overall')
axes[0].set_ylabel('Прогноз')
axes[0].set_title('Базовая модель: факт и прогноз')
axes[1].hist(residuals, bins=40)
axes[1].axvline(0, color='black', linewidth=1)
axes[1].set_xlabel('Остаток (факт − прогноз)')
axes[1].set_ylabel('Число игроков')
axes[1].set_title('Базовая модель: остатки на тесте')
plt.tight_layout()
plt.show()
\end{minted}

Точка на диагонали "--- точный прогноз. Остаток равен разности факта и прогноза. Сдвиг гистограммы от нуля показал бы систематическое завышение или занижение рейтинга.

\subsection{Подбор параметров и повторная оценка (пункты 5--7)}
\begin{minted}{python}
param_grid = {
    'rf__n_estimators': [100, 200],
    'rf__max_depth': [None, 12, 20],
    'rf__min_samples_leaf': [1, 4],
}

search = GridSearchCV(
    make_pipeline(),
    param_grid,
    cv=3,
    scoring='neg_root_mean_squared_error',
    n_jobs=1,
    refit=True,
)
search.fit(X_train, y_train)
print('Лучшие параметры:', search.best_params_)
print('RMSE на кросс-валидации:', round(-search.best_score_, 4))

best = search.best_estimator_
tuned_train = regression_metrics(y_train, best.predict(X_train))
tuned_test = regression_metrics(y_test, best.predict(X_test))
pd.DataFrame({
    'baseline train': baseline_train,
    'baseline test': baseline_test,
    'tuned train': tuned_train,
    'tuned test': tuned_test,
}).round(4)
\end{minted}

Сетка содержит 12 конфигураций и 3 фолда, то есть 36 обучений. Критерий "--- RMSE. \texttt{refit=True} заново учит лучший конвейер на всей обучающей части. Поиск идёт последовательно: внутри одного леса деревья уже считаются параллельно (\texttt{n\_jobs=-1}).

Для найденной модели строятся те же два графика, что и для базовой: факт против прогноза и гистограмма остатков на тесте.

\subsection{Важность признаков (пункты 8 и 9)}
\begin{minted}{python}
feature_names = best.named_steps['prep'].get_feature_names_out()
importance = pd.Series(
    best.named_steps['rf'].feature_importances_,
    index=[name.split('__', 1)[-1] for name in feature_names],
)
importance = importance.sort_values(ascending=True)

fig, ax = plt.subplots(figsize=(8, 6))
importance.plot.barh(ax=ax, color='steelblue')
ax.set_xlabel('Важность')
ax.set_title('Вклад признаков в случайный лес')
plt.tight_layout()
plt.show()

ranked = importance.sort_values(ascending=False)
print(ranked.head(5).round(4).to_string())
print('Доля топ-3:', round(float(ranked.head(3).sum()), 4))
\end{minted}

Имена берутся из матрицы, которую конвейер отдаёт лесу. Префиксы групп \texttt{money}, \texttt{num} и \texttt{cat} отбрасываются, на графике остаются исходные названия признаков и уровни индикаторов.

%======================== 6. ВЫВОДЫ ========================
\newpage
\section{Выводы по лабораторной работе}

В ходе работы на 18\,207 игроках FIFA обучен регрессионный случайный лес для рейтинга \texttt{Overall}. После отбора столбцов в таблице 14 признаков, часть из них ещё содержит пропуски: \texttt{Value} "--- 252, год контракта "--- 289, международная репутация и финты "--- по 48. Эти пропуски закрывает медиана или мода внутри конвейера, статистики не считаются по тесту и не используют цель. Обучающая часть содержит 14\,565 объектов, тестовая "--- 3\,642. Средний рейтинг на обучении $66{,}23$ при стандартном отклонении $6{,}93$, на тесте $66{,}29$ и $6{,}82$. Константный прогноз средним ошибся бы примерно на 7 пунктов. Корреляция с целью до очистки наибольшая у \texttt{Potential} ($0{,}661$), \texttt{Value} ($0{,}632$), \texttt{Release Clause} ($0{,}598$) и \texttt{Wage} ($0{,}572$).

Базовый лес "--- 100 деревьев, глубина не ограничена, лист из одного объекта. На обучении RMSE $0{,}176$, MAE $0{,}091$, $R^2 = 0{,}999$, MAPE $0{,}14~\%$. На тесте RMSE $0{,}478$, MAE $0{,}240$, MSE $0{,}229$, $R^2 = 0{,}995$, MAPE $0{,}38~\%$. Разрыв обучения и теста показывает, что глубокие деревья запоминают выборку плотнее, чем обобщают её. Сама тестовая ошибка при этом меньше половины пункта рейтинга и далека от ошибки константы. Точки «факт "--- прогноз» лежат вдоль диагонали, остатки собраны около нуля.

\texttt{GridSearchCV} по трём фолдам выбрал 200 деревьев, глубину 20 и лист из одного объекта. RMSE на кросс-валидации равен $0{,}507$. На тесте у этой модели MSE $0{,}224$, MAE $0{,}239$, RMSE $0{,}474$, $R^2 = 0{,}995$, MAPE $0{,}38~\%$. Сдвиг RMSE относительно базового леса "--- около $0{,}005$ пункта. Увеличение числа деревьев стабилизировало среднее, а ограничение глубины до 20 почти не изменило уже насыщенное качество. Минимальный лист сетка оставила равным одному: укрупнение листьев до четырёх объектов на этой зависимости не выигрывало у RMSE.

Почти весь вклад приходится на \texttt{Value}: важность $0{,}881$. Дальше идут \texttt{Age} ($0{,}076$) и \texttt{Potential} ($0{,}038$). Вместе три признака дают $0{,}994$ суммарной важности. \texttt{Release Clause} и \texttt{Weight} остаются на уровне $0{,}001$, позиция и рабочая нога на графике практически не видны. Такое распределение согласуется с устройством данных FIFA: стоимость игрока считается из рейтинга, поэтому \texttt{Value} почти повторяет цель, хотя линейная корреляция у потенциала выше. Лес берёт тот признак, по которому порог сразу отделяет уровни \texttt{Overall}, а не тот, у которого корреляция Пирсона максимальна. При известной стоимости возраст и потенциал лишь поправляют прогноз, а амплуа почти ничего не добавляет.

Высокий $R^2$ поэтому не стоит читать как открытие сложной зависимости. Модель честно использует столбец, в котором рейтинг уже зашит формулой игры. На признаках без стоимости ошибка была бы другой; сравнение базового леса и подобранного от этого не меняется: оба упираются в один и тот же потолок. Стандартизация и \texttt{log1p} на качество дерева почти не влияют, потому что разбиение монотонно инвариантно, и стоят в конвейере по условию работы.

Все пункты задания выполнены. Случайный лес предсказывает \texttt{Overall} со средней абсолютной ошибкой около $0{,}24$ пункта на тесте. Подбор числа деревьев и глубины улучшает RMSE незначительно. Основной вклад вносит стоимость игрока, затем возраст и потенциал.
\end{document}
